\chapter{Iodine Deficiency}

\section*{Definition and Overview}
Iodine deficiency is an inadequate intake of the mineral iodine, which the body needs to produce thyroid hormones. It is the world's most common preventable cause of intellectual disability, because thyroid hormones are essential for brain development in the fetus and infant. In adults, iodine deficiency causes goitre, an enlargement of the thyroid gland, and can progress to hypothyroidism. In Australia, mandatory iodisation of bread has largely corrected the mild deficiency that re-emerged in the 1990s, but pregnant and breastfeeding women, and people on restricted diets, remain at risk. Prevention through diet and supplementation is simple and highly effective.

\section*{Epidemiology}
Iodine deficiency affects around two billion people worldwide, making it a major global health problem, particularly in mountainous regions and inland areas where soil and food are iodine-poor. In Australia, bread iodisation introduced in 2009 restored population sufficiency, but mild deficiency persists in some groups. Pregnant and breastfeeding women have increased iodine needs and are the group most at risk of deficiency, followed by vegans and people avoiding iodised products. Severe deficiency in pregnancy causes cretinism and high rates of miscarriage and stillbirth in affected regions; in Australia, severe deficiency is essentially eliminated, but mild deficiency still warrants attention.

\section*{Aetiology and Pathophysiology}
The body needs around 150 micrograms of iodine daily for adults and more during pregnancy and breastfeeding. Iodine is absorbed from food and concentrated by the thyroid gland, which uses it to manufacture the hormones thyroxine (T4) and triiodothyronine (T3). When intake is inadequate, the thyroid enlarges to trap more iodine, producing a goitre, and as stores fall, hormone production drops, causing hypothyroidism. The most severe effects occur during fetal and early childhood development, when thyroid hormones direct the maturation of the brain: deficiency then causes permanent cognitive impairment, and in extreme cases cretinism, with intellectual disability, deafness, and growth failure. The thyroid is also vulnerable to the effects of deficiency through its increased drive for growth.

\section*{Clinical Features}
\begin{itemize}
    \item Goitre: a visible swelling at the base of the neck from an enlarged thyroid
    \item Symptoms of hypothyroidism: fatigue, weight gain, cold intolerance, dry skin, hair loss, and constipation
    \item In pregnancy: increased risk of miscarriage, stillbirth, and low birth weight
    \item In infants and children: delayed development, poor growth, and learning difficulties
    \item In severe deficiency: cretinism with intellectual disability and deafness
    \item Most mild deficiency produces no symptoms
    \item Deficiency may worsen the effects of other thyroid conditions
\end{itemize}

\section*{Differential Diagnosis}
Goitre and hypothyroidism have several causes beyond iodine deficiency.
\begin{itemize}
    \item \textbf{Hashimoto's thyroiditis:} the most common cause of hypothyroidism in Australia
    \item \textbf{Graves' disease:} autoimmune hyperthyroidism with goitre
    \item \textbf{Thyroid nodules and cancer:} lumps within the thyroid
    \item \textbf{Subacute thyroiditis:} transient painful thyroid inflammation
    \item \textbf{Medicines and dietary goitrogens:} substances that impair thyroid function
    \item \textbf{Congenital hypothyroidism:} detected by newborn screening
    \item \textbf{Other causes of fatigue and weight change:} depression, anaemia, and diabetes
\end{itemize}

\section*{When to Seek Medical Care}
People with a neck swelling, symptoms of an underactive thyroid, or dietary concerns should be assessed. Pregnant women and women planning pregnancy should discuss iodine supplementation with their doctor, and breastfeeding women should continue it. Anyone with known thyroid disease should have their iodine status considered before taking supplements.

\textbf{Call 000 (or local emergency services) for:}
\begin{itemize}
    \item Sudden difficulty breathing or swallowing from a rapidly enlarging goitre
    \item Rapid heart rate, fever, agitation, and confusion, which may indicate thyroid storm
    \item Severe neck pain or swelling
\end{itemize}

\section*{Diagnosis and Investigations}
\begin{itemize}
    \item \textbf{Thyroid function tests:} TSH, free T4, and free T3
    \item \textbf{Urinary iodine:} a measure of recent iodine intake
    \item \textbf{Thyroid antibodies:} to distinguish autoimmune disease from deficiency
    \item \textbf{Ultrasound:} assessment of goitre and nodules
    \item \textbf{Dietary assessment:} review of iodine-containing foods and supplements
    \item \textbf{Newborn screening:} congenital hypothyroidism screening in all Australian newborns
    \item \textbf{Population monitoring:} national surveys of urinary iodine
\end{itemize}

\section*{Management}
\begin{itemize}
    \item \textbf{Dietary correction:} increasing intake of iodised salt, seafood, and iodine-fortified bread
    \item \textbf{Supplementation:} iodine tablets for deficiency and in pregnancy and breastfeeding
    \item \textbf{Thyroxine treatment:} for established hypothyroidism
    \item \textbf{Monitoring:} repeat thyroid function and iodine tests to confirm correction
    \item \textbf{Pregnancy care:} supplementation from early pregnancy and through breastfeeding
    \item \textbf{Public health programs:} salt and bread iodisation to protect populations
    \item \textbf{Education:} advice for at-risk groups on dietary sources and supplements
\end{itemize}

\section*{Complications}
\begin{itemize}
    \item Goitre and compression of the airway or swallowing
    \item Hypothyroidism with its effects on metabolism, heart, and mood
    \item Miscarriage, stillbirth, and congenital abnormalities in pregnancy
    \item Permanent intellectual disability and cretinism in severe early deficiency
    \item Impaired growth and development in children
    \item Interaction with other thyroid disorders
\end{itemize}

\section*{Prognosis and Outcomes}
Iodine deficiency is easily prevented and treated. Population iodisation programs have transformed outcomes worldwide, eliminating cretinism where they are implemented, and Australia's bread iodisation restored population sufficiency. Established goitre and hypothyroidism respond well to thyroxine, and catch-up growth is possible in many children once intake is corrected. The crucial window is pregnancy and infancy, and supplementation at that time protects the developing brain. With adequate iodine, the complications of deficiency are almost entirely avoidable.

\section*{Prevention and Self-Care}
\begin{itemize}
    \item Use iodised salt and eat seafood, fish, eggs, and dairy
    \item Include iodine-fortified bread in the diet
    \item Take 150 micrograms of iodine daily in pregnancy and while breastfeeding
    \item Start supplementation when planning pregnancy
    \item Seek dietary advice if vegan or on a restricted diet
    \item Avoid excess iodine, including kelp and high-dose supplements
    \item Report neck swellings or symptoms of thyroid dysfunction
    \item Attend newborn and thyroid screening
\end{itemize}

\section*{Sources and Further Reading}
The following reputable sources provide further information for healthcare students and patients seeking reliable, current advice about iodine deficiency.
\begin{itemize}
    \item Thyroid Australia. \textit{Iodine deficiency information.} https://www.thyroid.org.au/
    \item National Health and Medical Research Council. \textit{Iodine supplementation in pregnancy.} https://www.nhmrc.gov.au/
    \item Food Standards Australia New Zealand. \textit{Iodine fortification.} https://www.foodstandards.gov.au/
    \item Healthdirect Australia. \textit{Iodine.} https://www.healthdirect.gov.au/
    \item World Health Organization. \textit{Iodine deficiency disorders.} https://www.who.int/
    \item UNICEF. \textit{Iodine deficiency global programs.} https://www.unicef.org/
\end{itemize}
